\documentclass{article}
\usepackage{array}
\usepackage[margin=.5in]{geometry}
\usepackage{titlesec}
\usepackage{longtable} % Add the longtable package
\usepackage{ulem}

\titleformat{\section}
  {\centering\footnotesize\bfseries} % format
  {\thesection} % label
  {1em} % horizontal space between label and title
  {\underline} % before-code

\begin{document}
\textbf{Name:}\underline{Name Goes Here} \hspace{90mm}\textbf{Date Of Lesson:}\underline{MM/DD/YYYY}

\begin{longtable}{ | m{50em} | } % Use longtable instead of tabular
  \hline
  \textbf{Content Area:} \\ 
  %fill in here- whenever you have \hline following what you write you must have \\ or it will compile poorly.
  \hline
  \textbf{Course Title and Grade Level:} \\ 
  \hline
  
  \underline{\textbf{Standards:}}\\ 
  %Standards:  Use the Nebraska Content Standards. Include 1-3 content standards.  List the specific standards for your content & grade—not just the number. Underline the portions of the standard that you will be addressing in this lesson. 
  \begin{center}
      \begin{itemize}
          \item 
          \item 
      \end{itemize}
  \end{center}\\
  \hline
  \underline{\textbf{Objectives:}}\\
  %Note: The content objective should be based on grade-level, content area standard(s). 
%Example: SWBAT state three of five central causes of the Civil War in writing. 
  \textbf{CO:}
  \vspace{.1in}
  
    %Note: You will say and/or write this objective so that all students have access to the learning goal. 
    %Examples: I can write about at least three causes of the Civil War OR Today I am learning more about the Civil War, so that I can state central causes of the war. I’ll know I’ve got it when I can write about at least three before the end of class.
  \textbf{SFO:} \\
  \hline
  \underline{\textbf{Assessments:}}\\
    %List the specific assessment strategies you are using and identify them as formative or summative. How will you determine if ALL students have met the above objectives?  It must be more than observation by walking around.  
   \begin{itemize}
        \item \footnotesize (1)
        \item \footnotesize (2)
        \item \footnotesize (3)
        \item \footnotesize (4)
        \item \footnotesize (5)
        \item \footnotesize (6)
   \end{itemize}\\
  \hline
  \underline{\textbf{Proactive Management:}}
    %Expectations (Voice-Movement-Task and materials), grouping strategy, productive discussions and small group roles. This is where you plan ahead to keep everything moving forward and to avoid any issues in your classroom! 

    \begin{itemize}
        \item \footnotesize (1)
        \item \footnotesize (2)
        \item \footnotesize (3)
        \item \footnotesize (4)
        \item \footnotesize (5)
        \item \footnotesize (6)
    \end{itemize}\\
  \hline 
  \underline{\textbf{Instructional Strategies:}}\\
    %List the specific instructional strategies (A MINUMUM OF 6 REQUIRED) you are using in this lesson. Discussion is not a strategy; which specific discussion strategy are you using? Then, describe in a few sentences how you will use them in your lesson plan.  
    %Name the strategy AND explain how you will use it and where!!!! 

  \begin{itemize}
        \item \footnotesize (1)
        \item \footnotesize (2)
        \item \footnotesize (3)
        \item \footnotesize (4)
        \item \footnotesize (5)
        \item \footnotesize (6)
    \end{itemize}\\
  \hline
  \underline{\textbf{Note Taking Strategies:}}\\
    %What specific strategy are students using?  How will they know what to write down/record?  

    %MUST INCLUDE FORMAL NOTE TAKING TEMPLATE/FORMAT FOR EACH LESSON. THIS MEANS, REGULAR, LINED PAPER DOES NOT WORK. Guided template that will encourage students to actually fill in all the “boxes/circles/etc” so that students are more likely to copy all the notes from class.  
  \footnotesize Here you type your note taking strategies\\
  \hline 
  \underline{\textbf{Questions:}}\\
    %List specific questions used throughout the lesson. Consider the following: Are your questions a blend of convergent and divergent?  What level of Bloom’s Taxonomy are your questions?  How will you ask and how will students respond (whole group, pairs, or small groups)? 

    %Question that TEACHER will ask students. 
  \textbf{Convergent Questions:}
  \begin{itemize}
        \item \footnotesize (1)
        \item \footnotesize (2)
        \item \footnotesize (3)
        \item \footnotesize (4)
        \item \footnotesize (5)
        \item \footnotesize (6)
    \end{itemize}
    \textbf{Divergent Questions:}
    \begin{itemize}
        \item \footnotesize (1)
        \item \footnotesize (2)
        \item \footnotesize (3)
        \item \footnotesize (4)
        \item \footnotesize (5)
        \item \footnotesize (6)
    \end{itemize}\\
  \hline 
  \begin{center}
      \underline{\textbf{Beginning Of The Lesson:}}
  \end{center}
% This may include bellwork and/or an anticipatory set to launch the lesson. Review: Activate Prior Knowledge Preview: State the objective and where the lesson is going Hook: Motivate and engage ALL learners with relevant connections. 
  \textbf{Bell work:}\\
  \underline{\textbf{Anticipatory Set:}}\\
  \begin{center}
    \textbf{Review:}\\
    \textbf{Preview:}\\
    \textbf{Hook:}\\  
    \end{center}
    \section*{Minute 0-3}
\footnotesize \textbf{Pacing:} \\
\footnotesize \textbf{Transition:} \\
\footnotesize \textbf{Content(Problem)}\\
\footnotesize \textbf{Directions/Activity}\\
  \hline 
  \begin{center}
      \underline{\textbf{The Lesson:}}
  \end{center}
%Outline of procedures including content, UDL, transitions, directions, pacing (provide estimated time), questions, check for learning, etc. Think of a detailed “sub-plan” that is easy to read and reference (e.g., avoid blocks of text/paragraphs) No more than 6 minutes intervals here! Explicitly label Gradual Release of Responsibility components MUST CONTAIN I-DO YOU-DO WE-DO
\section*{Minute 3-6}
\footnotesize \textbf{Pacing:} \\
\footnotesize \textbf{Transition:} \\
\footnotesize \textbf{Content(Problem)}\\
\footnotesize (1)\\
\footnotesize (2) \\
\footnotesize (3)\\
\footnotesize \textbf{Directions/Activity}\\
\footnotesize (1) \\
\footnotesize (2) \\

\section*{Minute 6-10}

\footnotesize \textbf{Pacing:} \\
\footnotesize \textbf{Transition:} \\
\footnotesize \textbf{Content(Problem)}\\
\footnotesize \textbf{Directions/Activity}\\

  \hline 
  \begin{center}
      \underline{\textbf{Closure:}}
  \end{center}\\
%What will you do or say to close the lesson? 
%Procedural— What do the students (or you the teacher) need to do in your room so that they can leave? (E.g., pick up scissors, assign homework, return quiz, etc.) 
%Content Summary— Circle back to objective.  You will say this summary to students! 
  \textbf{Content Summary:}\\
  \textbf{Procedural:}\\
  \hline
  \textbf{References:}\\
%Cite any texts, websites, videos, etc. used during your lesson in APA format. If you are teaching from your mentor’s lesson plan, please state here to allocate proper planning credit. You should include at least 6 references here.
  \hline 
\end{longtable} % Close the longtable environment

\section*{Lesson Plan Reflection:}

\textbf{Highlight Culturally Responsive Teaching (CRT) components within your lesson.} 

\textbf{Italicize Universal Design for Learning (UDL) components within your lesson. }

\begin{enumerate}
    \item Write a paragraph (3-5 sentences) elaborating on the intentionally planned CRT components of your lesson. How did you leverage student assets, thinking, language and cultural backgrounds to make learning relevant while demonstrating high expectations for all students? 
    \item Write a paragraph (3-5 sentences) elaborating on the intentionally planned UDL components (representation, action and expression, engagement) of your lesson. What strategies did you use to meet the needs of ALL learners? 

\end{enumerate}

\end{document}
